\documentclass[10pt,a4paper]{amsart}
\usepackage{amsmath,amssymb,amsthm,graphicx}
\usepackage{verbatim}
\newenvironment{reference}{\comment}{\endcomment}
%\newenvironment{reference}{}{}
\newenvironment{slogan}{\comment}{\endcomment}
\newenvironment{history}{\comment}{\endcomment}

% For commutative diagrams we use Xy-pic
\usepackage[all]{xy}

% We use 2cell for 2-commutative diagrams.
\xyoption{2cell}
% Theorem environments.
%
\theoremstyle{plain}
\newtheorem{theorem}[subsection]{Theorem}
\newtheorem{proposition}[subsection]{Proposition}
\newtheorem{lemma}[subsection]{Lemma}

\theoremstyle{definition}
\newtheorem{definition}[subsection]{Definition}
\newtheorem{example}[subsection]{Example}
\newtheorem{exercise}[subsection]{Exercise}
\newtheorem{situation}[subsection]{Situation}

\theoremstyle{remark}
\newtheorem{remark}[subsection]{Remark}
\newtheorem{remarks}[subsection]{Remarks}

\numberwithin{equation}{subsection}

% Macros
%
\def\lim{\mathop{\mathrm{lim}}\nolimits}
\def\colim{\mathop{\mathrm{colim}}\nolimits}
\def\Spec{\mathop{\mathrm{Spec}}}
\def\Hom{\mathop{\mathrm{Hom}}\nolimits}
\def\Ext{\mathop{\mathrm{Ext}}\nolimits}
\def\SheafHom{\mathop{\mathcal{H}\!\mathit{om}}\nolimits}
\def\SheafExt{\mathop{\mathcal{E}\!\mathit{xt}}\nolimits}
\def\Sch{\mathit{Sch}}
\def\Mor{\mathop{\mathrm{Mor}}\nolimits}
\def\Ob{\mathop{\mathrm{Ob}}\nolimits}
\def\Sh{\mathop{\mathit{Sh}}\nolimits}
\def\NL{\mathop{N\!L}\nolimits}
\def\CH{\mathop{\mathrm{CH}}\nolimits}
\def\proetale{{pro\text{-}\acute{e}tale}}
\def\etale{{\acute{e}tale}}
\def\QCoh{\mathit{QCoh}}
\def\Ker{\mathop{\mathrm{Ker}}}
\def\Im{\mathop{\mathrm{Im}}}
\def\Coker{\mathop{\mathrm{Coker}}}
\def\Coim{\mathop{\mathrm{Coim}}}

% Boxtimes
%
\DeclareMathSymbol{\boxtimes}{\mathbin}{AMSa}{"02}

%
% Macros for moduli stacks/spaces
%
\def\QCohstack{\mathcal{QC}\!\mathit{oh}}
\def\Cohstack{\mathcal{C}\!\mathit{oh}}
\def\Spacesstack{\mathcal{S}\!\mathit{paces}}
\def\Quotfunctor{\mathrm{Quot}}
\def\Hilbfunctor{\mathrm{Hilb}}
\def\Curvesstack{\mathcal{C}\!\mathit{urves}}
\def\Polarizedstack{\mathcal{P}\!\mathit{olarized}}
\def\Complexesstack{\mathcal{C}\!\mathit{omplexes}}
% \Pic is the operator that assigns to X its picard group, usage \Pic(X)
% \Picardstack_{X/B} denotes the Picard stack of X over B
% \Picardfunctor_{X/B} denotes the Picard functor of X over B
\def\Pic{\mathop{\mathrm{Pic}}\nolimits}
\def\Picardstack{\mathcal{P}\!\mathit{ic}}
\def\Picardfunctor{\mathrm{Pic}}
\def\Deformationcategory{\mathcal{D}\!\mathit{ef}}
\usepackage{cleveref}
\crefname{originalreference}{original reference}{original references}
\Crefname{originalreference}{Original reference}{Original references}
\makeatletter
\newlabel{lemma-vanishing}{{07LA}{source}{}{sourcepointer-1}{}}
\newlabel{lemma-vanishing@cref}{{[originalreference][0][]07LA}{[1][0][]source}{}{}{}}
\makeatother
\crefname{originalreference}{original reference}{original references}
\Crefname{originalreference}{Original reference}{Original references}
\makeatletter
\newlabel{situation-Cohen-Macaulay-curve}{{0B5D}{source}{}{sourcepointer-2}{}}
\newlabel{situation-Cohen-Macaulay-curve@cref}{{[originalreference][0][]0B5D}{[1][0][]source}{}{}{}}
\makeatother
\title{Surjectivity of the globally generated submodule onto a zero-dimensional quotient}
\author{The Stacks Project Authors}
\date{stacks-project:curves.tex:local-snapshot-20261012; GNU FDL 1.2 or later}
\begin{document}
\maketitle
\noindent\textit{Editorial note.} All delivered mathematical text is the author's own, copied from the pinned Stacks Project chapter file sources/stacks-curves/pinned-original-source.tex; no statement or proof step was rewritten, reordered or removed.\par\medskip
\noindent\textit{Editorial note.} The counted proof is the author's complete argument: the closing contradiction 'H\textasciicircum{}0(X, F) is bigger than H\textasciicircum{}0(X, F')' (source lines 1399-1401) is resolved by the definition of F' as the submodule generated by ALL global sections of F: any global section of F factors through F', so Gamma(F) = Gamma(F') and the quotient H\textasciicircum{}0(F)/H\textasciicircum{}0(F') is zero; no external input is needed for this step.\par\medskip
\noindent\textit{Editorial note.} The label lemma-globally-generated names the counted statement, and the proof environment immediately following it in the source is delivered verbatim as its proof.\par\medskip
\noindent\textit{Editorial note.} The counted statement is carried by the author's own lemma environment in the chapter source.\par\medskip
\noindent\textit{Allowed standard prerequisites (editorial).}\begin{itemize}
\item Cohomological vanishing for the constructed submodule. Same chapter, \textbackslash{}label\{lemma-vanishing\}. Proof lines 1391-1393: H\textasciicircum{}1(X, F') = 0 follows from this lemma and the assumption that no nonzero submodule of F injects into Q.
\item The Cohen-Macaulay curve situation. Same chapter, Situation \textbackslash{}label\{situation-Cohen-Macaulay-curve\}. Statement line 1372: the lemma is set in this situation; in particular dim\_k H\textasciicircum{}1(X, omega\_X) = dim\_k H\textasciicircum{}0(X, O\_X) = 1 (used in proof lines 1387-1389).
\end{itemize}\medskip
\ifdefined\hypertarget\hypertarget{sourcepointer-1}{}\fi
\noindent\textit{Source reference.} Original reference 07LA: sources/stacks-curves/pinned-original-source.tex, source lines 1321--1321. Prerequisite pointer invoked at proof line(s) 1392; it names a standard fact of the same chapter that the counted argument uses, and every step of the counted proof built on it is retained verbatim in the delivered body.\par\medskip
\ifdefined\hypertarget\hypertarget{sourcepointer-2}{}\fi
\noindent\textit{Source reference.} Original reference 0B5D: sources/stacks-curves/pinned-original-source.tex, source lines 1307--1307. Prerequisite pointer invoked at proof line(s) 1373; it names a standard fact of the same chapter that the counted argument uses, and every step of the counted proof built on it is retained verbatim in the delivered body.\par\medskip
\begin{lemma}
\label{lemma-globally-generated}
In Situation \ref{situation-Cohen-Macaulay-curve}. Given an exact sequence
$$
\omega_X \to \mathcal{F} \to \mathcal{Q} \to 0
$$
of coherent $\mathcal{O}_X$-modules with $\dim(\text{Supp}(\mathcal{Q})) = 0$
and $\dim_k H^0(X, \mathcal{Q}) \geq 2$ and such that there is no nonzero
submodule $\mathcal{Q}' \subset \mathcal{F}$ such that
$\mathcal{Q}' \to \mathcal{Q}$ is injective.
Then the submodule of $\mathcal{F}$ generated by global
sections surjects onto $\mathcal{Q}$.
\end{lemma}

\begin{proof}
Let $\mathcal{F}' \subset \mathcal{F}$ be the submodule generated by
global sections and the image of $\omega_X \to \mathcal{F}$. Since
$\dim_k H^0(X, \mathcal{Q}) \geq 2$ and
$\dim_k H^1(X, \omega_X) = \dim_k H^0(X, \mathcal{O}_X) = 1$,
we see that $\mathcal{F}' \to \mathcal{Q}$ is not zero and
$\omega_X \to \mathcal{F}'$ is not an isomorphism.
Hence $H^1(X, \mathcal{F}') = 0$ by Lemma \ref{lemma-vanishing}
and our assumption on $\mathcal{F}$.
Consider the short exact sequence
$$
0 \to \mathcal{F}' \to \mathcal{F} \to
\mathcal{Q}/\Im(\mathcal{F}' \to \mathcal{Q}) \to 0
$$
If the quotient on the right is nonzero, then we obtain a contradiction
because then $H^0(X, \mathcal{F})$ is bigger than $H^0(X, \mathcal{F}')$.
\end{proof}

\end{document}
